\chapter{ФИЗИЧЕСКАЯ МОДЕЛЬ И ЧИСЛЕННАЯ МЕТОДИКА}

\section{Гамильтониан квадратной решётки в методе сильной связи}

В работе рассматривается модельная квантовая точка, заданная как конечная область квадратной решётки. Такая постановка не претендует на описание конкретного материала на уровне первопринципного расчёта. Её назначение состоит в том, чтобы контролируемо исследовать влияние размера и формы конечной низкоразмерной наноструктуры на низкоэнергетический спектр. Квадратная решётка в этой модели играет роль простой кристаллической дискретности, а конечная область решётки соответствует модельной квантовой точке или конечному нанокристаллиту с заданной границей.

Для описания электронного спектра используется гамильтониан метода сильной связи с ближайшими соседями~\cite{slater1954lcao,ashcroft1976solidstate}:
\begin{equation}
  H
  =
  \sum_i \varepsilon_i |i\rangle\langle i|
  +
  \sum_{\langle i,j\rangle}
  t_{ij}
  \left(
    |i\rangle\langle j|
    +
    |j\rangle\langle i|
  \right),
\end{equation}
где $|i\rangle$ обозначает состояние, локализованное на узле решётки $i$, $\varepsilon_i$ --- энергия на узле, а $t_{ij}$ --- интеграл перескока между ближайшими соседями. Сумма $\langle i,j\rangle$ берётся по парам соседних узлов квадратной решётки.

В данной работе используется однородная безразмерная конвенция
\begin{equation}
  \varepsilon_i = 0,
  \qquad
  t_{ij} = -1
  \quad
  \text{для ближайших соседей}.
\end{equation}
Открытая граница возникает естественно: в гамильтониан включаются только узлы, принадлежащие выбранной геометрической области, а перескоки к узлам вне этой области отсутствуют. Поэтому граница модельной квантовой точки является жёсткой стенкой в рамках дискретной решёточной модели.

Такой гамильтониан задаёт прозрачную физическую модель квантового ограничения. Изменение размера области меняет характерные волновые числа низших состояний, а изменение формы границы влияет на симметрию и относительное положение уровней. Эти свойства являются центральными для дальнейшего анализа суперэллиптических квантовых точек.

\section{Энергетическая шкала и низкоэнергетический предел}

Для бесконечной квадратной решётки с нулевой энергией на узле и интегралом перескока $-1$ дисперсионное соотношение имеет вид
\begin{equation}
  E(k_x,k_y)
  =
  -2\cos k_x
  -
  2\cos k_y .
\end{equation}
Минимум этой зоны находится в точке $k_x=k_y=0$ и равен
\begin{equation}
  E_{\min} = -4 .
\end{equation}
При малых $k_x$ и $k_y$ разложение косинусов даёт квадратичный низкоэнергетический предел:
\begin{equation}
  E(k_x,k_y) + 4
  \approx
  k_x^2 + k_y^2 .
\end{equation}

Поэтому для сравнения низших уровней с континуальными оценками используется энергия, отсчитанная от дна зоны:
\begin{equation}
  E_{\mathrm{kin}} = E_0 + 4 .
\end{equation}
Здесь $E_0$ --- нижнее собственное значение конечной системы в выбранной конвенции модели сильной связи. Сама величина $E_0$ содержит сдвиг, связанный с положением дна зоны на квадратной решётке, тогда как $E_{\mathrm{kin}}$ характеризует кинетическую энергию низкоэнергетического состояния относительно этого дна.

Эта энергетическая шкала важна для дальнейшей физической интерпретации. Если конечная система находится в низкоэнергетическом квадратичном режиме, то характерное волновое число ограниченного состояния масштабируется как обратный размер области, а энергия $E_{\mathrm{kin}}$ должна иметь ведущую зависимость порядка $1/a^2$. Подробная численная проверка этого масштабирования приведена в главе 3.

\section{Геометрия конечной квантовой точки}

Конечная система строится как множество узлов квадратной решётки, координаты которых удовлетворяют заданному геометрическому условию. Такая конструкция переводит гладкую границу модельной наноструктуры в дискретную границу решётки. Именно эта дискретная граница затем влияет на спектр наряду с общим размером и формой области.

Прямоугольная область используется как контрольная задача. Для прямоугольника с размерами $N_x \times N_y$ и открытыми границами спектр квадратной решётки разделяется по двум координатам:
\begin{equation}
\begin{aligned}
  E_{mn}
  &=
  -2
  \left[
    \cos\left(\frac{\pi m}{N_x+1}\right)
    +
    \cos\left(\frac{\pi n}{N_y+1}\right)
  \right], \\
  &\quad
  m=1,\ldots,N_x,\quad n=1,\ldots,N_y .
\end{aligned}
\end{equation}
Эта формула используется не как основной объект исследования, а как проверка энергетической конвенции, открытых границ и численной процедуры.

Основное семейство геометрий в работе задаётся суперэллипсом:
\begin{equation}
  \left|\frac{x}{a}\right|^n
  +
  \left|\frac{y}{b}\right|^n
  \leq 1,
  \qquad
  b = a r_{\mathrm{AR}} .
\end{equation}
Параметр $a$ задаёт характерный размер вдоль оси $x$, параметр $b$ --- размер вдоль оси $y$, а $r_{\mathrm{AR}}$ является отношением сторон. Показатель $n$ управляет формой границы. При $n=2$ получается эллиптическая, а при $r_{\mathrm{AR}}=1$ круговая форма. При больших $n$ граница становится более прямоугольноподобной, а при $n$, близком к единице, --- более ромбоподобной; при рассматриваемых $n\geq 1$ область остаётся выпуклой.

В дальнейших расчётах используются фиксированные классы формы
\[
  n = 1.2,\;2.0,\;3.0,\;4.0 .
\]
Параметр $n$ не рассматривается как непрерывный входной параметр суррогатной модели. Поэтому результаты работы относятся к указанным дискретным классам формы и не являются утверждением об обобщении на произвольные значения $n$ или произвольные границы.

\section{Численный расчёт спектра в Kwant}

Для построения конечных решёточных систем и вычисления их спектров используется Kwant~\cite{groth2014kwant}. Физическая постановка при этом задаётся гамильтонианом метода сильной связи, а Kwant выполняет численную реализацию этой постановки: создаёт конечную систему, присваивает энергии на узлах и интегралы перескока, формирует матрицу гамильтониана и передаёт её алгоритму поиска собственных значений.

Расчётная процедура состоит из нескольких шагов. Сначала задаётся геометрическая область: прямоугольник или суперэллипс. Затем выбираются все узлы квадратной решётки, лежащие внутри этой области. Для каждого такого узла задаётся энергия на узле $\varepsilon_i=0$, а между ближайшими соседями внутри области задаётся интеграл перескока $t_{ij}=-1$. Связи с узлами, не входящими в область, не добавляются, что соответствует открытой границе конечной системы.

После построения конечной системы получается разреженная эрмитова матрица гамильтониана. Для данной работы требуются только нижние уровни спектра, поэтому численно вычисляются несколько наименьших собственных значений; для таких задач стандартны итерационные методы поиска собственных значений разреженных матриц~\cite{lehoucq1998arpack}. Эти значения являются прямым спектральным результатом модели и используются как эталон для всех последующих суррогатных аппроксимаций. Суррогатные модели, рассматриваемые в следующих главах, не заменяют этот расчёт, а только приближают уже полученные спектральные зависимости внутри проверенной области параметров.

\section{Целевые спектральные характеристики}

Пусть собственные значения конечного гамильтониана упорядочены по возрастанию:
\begin{equation}
  E_0 \leq E_1 \leq E_2 \leq E_3 \leq \ldots .
\end{equation}
Величина $E_0$ является нижним уровнем конечной модельной квантовой точки. Она характеризует основной уровень квантового ограничения в выбранной геометрии. Первый зазор определяется как
\begin{equation}
  dE_1 = E_1 - E_0 .
\end{equation}
Он описывает расстояние от основного уровня до первого возбуждённого состояния и поэтому чувствителен как к размеру, так и к форме границы.

Следующий зазор можно записать как
\begin{equation}
  dE_2 = E_2 - E_1 .
\end{equation}
Однако эта величина не используется как основная целевая величина суррогатной регрессии. В симметричных геометриях соседние возбуждённые уровни могут быть вырождены или почти вырождены, и тогда $dE_2$ становится чувствительным к порядку уровней и малым решёточным деталям границы. Поэтому в данной работе основными спектральными характеристиками являются $E_0$ и $dE_1$, а $dE_2$ используется как диагностическая величина. Подробная численная иллюстрация этого решения приведена в главе 3.

\section{Выводы по главе 2}

В данной главе задана физическая и численная основа дальнейшего исследования. Основные положения состоят в следующем.

\begin{enumerate}
  \item Модельная квантовая точка описывается гамильтонианом метода сильной связи на квадратной решётке с нулевой энергией на узле и интегралом перескока $-1$ между ближайшими соседями.
  \item Нижний край зоны бесконечной квадратной решётки равен $-4$, поэтому для низкоэнергетической интерпретации используется величина $E_{\mathrm{kin}}=E_0+4$.
  \item Конечные области решётки моделируют квантовые точки или конечные нанокристаллиты с открытой границей; прямоугольник служит контрольной задачей, а суперэллипс является основной геометрической семьёй.
  \item Параметры $a$, $r_{\mathrm{AR}}$ и $n$ управляют размером, отношением сторон и формой границы, что определяет спектр через квантовое ограничение.
  \item Kwant используется как инструмент прямого численного расчёта спектра конечного гамильтониана метода сильной связи; эти расчёты являются эталонными для последующих суррогатных моделей.
  \item Основными спектральными характеристиками выбраны $E_0$ и $dE_1$, тогда как $dE_2$ рассматривается как диагностическая величина из-за возможного вырождения и чувствительности порядка уровней.
\end{enumerate}

\clearpage
